\section{Tìm kiếm theo chiều rộng (Breadth-First Search, BFS)}
\label{sec:bloxorz-bfs}

BFS sử dụng chung \texttt{State}, \texttt{Action}, \texttt{Board} và vòng lặp tìm kiếm
chung của solver cơ sở. Sự khác biệt của BFS nằm ở chính sách lấy các node: hàng đợi
FIFO.

\subsection{Thứ tự FIFO và chiến lược}

Đối với BFS, frontier là hàng đợi FIFO. Hai phương thức quan trọng trong hàng đợi gồm
\texttt{push} và \texttt{pop}. Trong đó, \texttt{push} sẽ không duyệt lại các node đã
gặp và \texttt{pop} sẽ lấy phần tử cũ nhất ra khỏi frontier. Các node được đánh dấu là
đã gặp tại thời điểm sinh node, không phải lúc mở rộng frontier. Do đó mà mỗi trạng thái
chỉ vào frontier đúng một lần nên vòng lặp tìm kiếm không phụ thuộc vào chu trình.

Kiểm tra đích diễn ra khi lấy nút ra, và các nút được lấy theo độ sâu không giảm. Nút
đích đầu tiên được lấy ra có độ sâu nhỏ nhất; vì BFS không đọc chi phí, mỗi cạnh trong
cây tìm kiếm đúng bằng một bước, nên độ sâu chính là số bước. Do đó BFS bảo đảm tính đầy
đủ trên đồ thị hữu hạn và trả về đường ít bước nhất. \texttt{bfs.py} chỉ gọi lại
\texttt{search(stage, caps, frontier=Frontier())}: luật trò chơi nằm ở
\texttt{Board.successors}, còn \texttt{Frontier()} đã là hàng đợi FIFO.

\subsection{Quan hệ với mô hình chi phí}

BFS dùng FIFO và không đọc \texttt{step\_cost}, nên thuật toán chỉ tối ưu số bước di
chuyển, không tối ưu chi phí. Trên các màn chỉ gồm ô thường, mọi bước đi đều có giá trị
1, nên số bước trùng với chi phí và BFS tối ưu cả hai. Trên màn có ô dễ vỡ, chi phí mỗi
bước không còn bằng 1 nên hai mục tiêu tách nhau: BFS có thể trả về đường ít bước hơn
nhưng đắt hơn UCS (ví dụ ở mục~\ref{sec:bloxorz-ucs}).

\subsection{Vòng lặp chung, bộ đếm và giới hạn}

Phần cài đặt của BFS gọi \texttt{base.search(stage, caps, frontier=Frontier())}, khai
báo frontier tường minh và đối xứng với \texttt{ucs.py} và \texttt{astar.py}, dù
\texttt{Frontier()} vốn đã là FIFO mặc định. Theo quy ước của \#18, \texttt{generated}
đếm cả các node bị frontier từ chối, \texttt{expanded} tăng trước \texttt{successors};
\texttt{frontier\_peak} là số node sống.

Giới hạn thời gian được kiểm tra trước khi lấy nút; giới hạn số nút nằm sau kiểm tra
đích và trước khi tăng \texttt{expanded}. Mọi kết quả giải được, đạt giới hạn hay không
có nghiệm đều trả \texttt{Counts}. Bộ giải chỉ dùng đồng hồ để thực thi giới hạn; thời
gian và bộ nhớ thực nghiệm do công cụ đo chung phụ trách, nên mục này không nêu số đo.
Gọi $V$ và $E$ là số trạng thái và cạnh đạt được: BFS chạy trong $O(V+E)$ thời gian,
dùng $O(V)$ bộ nhớ cho tập trạng thái đã gặp và frontier, và không có hàng đợi ưu tiên
nên không phát sinh thừa số $\log$.

\subsection{Kiểm chứng và kết quả}

Ba màn có đáp án tính tay cho kết quả 6, 4 và 4 bước, trùng với
\texttt{expected.optimal\_cost}. Với mỗi màn, đường đi trả về được phát lại trên
\texttt{Board}; trạng thái cuối là đích và chi phí phát lại khớp chi phí bộ giải báo.
Trên mười màn chính thức, số bước của BFS được đối chiếu với một phép duyệt theo chiều
rộng viết độc lập ngay trong kiểm thử, không đi qua \texttt{base.search}; hai kết quả
khớp trên cả mười màn.

Hai kiểm thử nhỏ khoá hành vi của frontier. Màn thứ nhất đặt đích là nút con đầu tiên và
một nhánh năm bước là nút con thứ hai: FIFO lấy đích ngay và trả một bước, trong khi
LIFO sẽ đi hết nhánh và trả sáu bước. Màn thứ hai chứa chu trình $A\to B\to A$: vì trạng
thái được đánh dấu ngay khi sinh, $A$ bị từ chối ở lần gặp thứ hai, vòng lặp dừng và mỗi
trạng thái chỉ được mở rộng một lần, với ba nút mở rộng và năm nút sinh.

So với UCS, BFS cho cùng số nút mở rộng, số nút sinh và đỉnh frontier trên chín trong
mười màn. Ở các màn đó đường ít bước nhất cũng là đường rẻ nhất, nên hai chiến lược duyệt
cùng một tập nút dù BFS không đọc chi phí. Màn L2-02 là ngoại lệ: BFS mở rộng 134 nút,
sinh 286 nút, đỉnh frontier 9; UCS là 135, 288 và 10. BFS duyệt ít hơn nhưng trả đường
đắt hơn, vì mục tiêu số bước và mục tiêu chi phí đã tách nhau từ khi có ô dễ vỡ
(Mục~\ref{sec:bloxorz-ucs}). Đây là kiểm tra tính đúng, không phải số đo thời gian hay bộ
nhớ; chưa thể kết luận BFS nhanh hơn hay tốn ít bộ nhớ hơn UCS trước khi có số đo của
công cụ chung.
